\begin{exercise}{51}
Show the following assertions:
\begin{enumerate}[label=(\arabic*)]
\item Let $\mathcal C$ be a nonempty ordinary category\footnote{The original says ``an ordinary category'', without the nonempty hypothesis. Add ``nonempty'' in (1): for $\mathcal C=\varnothing$ and $X=\partial\Delta^2$, there are no maps $X\to\mathrm{N}(\mathcal C)=\varnothing$, so the condition is vacuous. Yet the inner horn with successive edges $0\to1\to2$ has no filler in $\partial\Delta^2$.} and $X$ a simplicial set. Then $X$ is an $\infty$-category if and
only if every map $X\to\mathrm{N}(\mathcal C)$ is an inner fibration.
\item A map $f\colon X\to Y$ is an inner fibration if and only if $f^{\mathrm{op}}\colon X^{\mathrm{op}}\to Y^{\mathrm{op}}$ is an inner
fibration.
\item A map $f\colon X\to Y$ is a left fibration if and only if the map $f^{\mathrm{op}}\colon X^{\mathrm{op}}\to Y^{\mathrm{op}}$ is a right
fibration.
\end{enumerate}

\begin{proof}
(1) If $X$ is an $\infty$-category, consider an inner-horn
lifting problem for $f\colon X\to\mathrm{N}(\mathcal C)$.
Fill its top horn in $X$. The image of this filler and the
prescribed bottom simplex are extensions of the same horn in
$\mathrm{N}(\mathcal C)$, so they agree by the uniqueness
in Theorem~1.1.52. This is a lift.
Conversely, choose an object $c\in\mathcal C$.
The constant map $X\to\mathrm{N}(\mathcal C)$ at $c$ is
an inner fibration by hypothesis.
Every inner horn in $X$, together with the constant simplex
at $c$, gives a lifting problem for this map.
Its lift fills the horn, proving that $X$ is an $\infty$-category.

(2)--(3) As in Exercise~\ref{ex:48}, taking opposites
identifies $\Lambda^n_j\hookrightarrow\Delta^n$ with
$\Lambda^n_{n-j}\hookrightarrow\Delta^n$.
Applied to an entire commutative lifting square and its diagonal,
this gives a bijection between lifting problems and their
solutions for $f$ and $f^{\mathrm{op}}$ at these respective indices.
The involution $j\mapsto n-j$ preserves the inner indices
$0<j<n$ and interchanges the left indices $0\le j<n$
with the right indices $0<j\le n$.
Both assertions now follow from Definition~1.3.1.
\end{proof}
\end{exercise}
